\documentclass[11pt]{article}

\usepackage[T1]{fontenc}
\usepackage{amsmath,amssymb,amsthm}
\usepackage{hyperref}

\title{Nested FIML Satorra--2000 in Saturated Eta Space}
\author{}
\date{\today}

\newcommand{\E}{\mathbb{E}}
\newcommand{\tr}{\operatorname{tr}}
\newcommand{\Var}{\operatorname{Var}}
\newcommand{\Cov}{\operatorname{Cov}}
\newcommand{\vech}{\operatorname{vech}}
\newcommand{\argmin}{\operatorname*{arg\,min}}
\newcommand{\argmax}{\operatorname*{arg\,max}}
\newcommand{\plim}{\operatorname*{plim}}
\newcommand{\rank}{\operatorname{rank}}
\newcommand{\eig}{\operatorname{eig}}
\newcommand{\T}{^{\mathsf T}}
\newcommand{\op}{o_p}
\newcommand{\Op}{O_p}

\theoremstyle{plain}
\newtheorem{proposition}{Proposition}
\newtheorem{corollary}{Corollary}
\theoremstyle{remark}
\newtheorem*{remark}{Remark}

\begin{document}
\maketitle

\section*{Purpose}

The question is whether magmaan's missing-data nested
\texttt{method = "restriction\_map"} statistic has an analytical target, or
whether it is only a diagnostic convention. The answer is conditional.

Under the ordinary nested null, where the smaller model is correct in the
saturated first-stage eta target, the eta-space statistic is the standard local
Satorra restricted-test derivation with the complete-data sample moments
replaced by the saturated FIML eta estimator. The finite-sample objects used by
magmaan,
\[
  V=\widehat H_\eta,\qquad
  \widehat\Gamma_{\mathrm{mis}}
  =\widehat H_\eta^{-1}\widehat J_\eta\widehat H_\eta^{-1},
\]
are plug-in estimates of the curvature and first-stage sandwich covariance.
The estimated curvature or two-stage weight does not add a first-order term at
this null, because both fitted residuals are \(\Op(N^{-1/2})\).
Thus ``robust'' here means robust to the data distribution and missing-data
first-stage covariance under the true nested null; it does not mean robust to
structural misspecification of the larger model.

The same derivation does not prove the statistic for the broader null in which
the nested restriction holds only at the pseudo-true point while the larger
model remains misspecified. In that regime the residual is nonzero, and the
limiting law is the Hessian of the two profiled discrepancies, including
curvature and estimated-weight derivative channels.

\section{Generic eta-space nested result}
\label{sec:generic}

Let \(\hat\eta\) be a first-stage estimator of a saturated moment vector
\(\eta_0\), with
\[
  Z_N=\sqrt N(\hat\eta-\eta_0)
  \Rightarrow Z,\qquad Z\sim N(0,\Gamma).
\]
Let the larger model \(M_1\) contain \(\eta_0\) and have local coordinates
\(\alpha\in\mathbb R^r\),
\[
  \eta_1(\alpha_N)
  =\eta_0+N^{-1/2}\Delta\alpha+\op(N^{-1/2}),
  \qquad \Delta\in\mathbb R^{q\times r}.
\]
The smaller model \(M_0\subset M_1\) is represented locally by
\[
  R\alpha=0,\qquad R\in\mathbb R^{m\times r},\qquad \rank(R)=m.
\]
Assume the local contrast for model \(M\) has the quadratic representation
\[
  T_M
  =
  \min_t (Z_N-\Delta_M t)\T V(Z_N-\Delta_M t)+\op(1),
  \label{eq:local-contrast}
\]
where \(V\) is symmetric positive definite and \(\Delta_M\) is the local tangent
matrix for model \(M\). For the larger model write
\[
  A=\Delta\T V\Delta,\qquad
  B=\Delta\T V\Gamma V\Delta.
\]

\begin{proposition}[Nested eta-space restriction map]
\label{prop:restriction}
Under the assumptions above,
\[
  T_0-T_1
  \Rightarrow
  (R\hat\alpha)\T C^{-1}(R\hat\alpha),
  \qquad
  \hat\alpha=A^{-1}\Delta\T VZ,
  \qquad
  C=RA^{-1}R\T.
  \label{eq:wald-form}
\]
Equivalently,
\[
  T_0-T_1
  \Rightarrow
  \sum_{j=1}^m\lambda_j\chi^2_{1j},
  \qquad
  \lambda_j
  =
  \eig_j(C^{-1}S),
  \label{eq:spectrum}
\]
where
\[
  S=RA^{-1}BA^{-1}R\T .
\]
\end{proposition}

\begin{proof}
The unconstrained minimizer in the larger tangent space is
\[
  \hat\alpha_N
  =
  A^{-1}\Delta\T VZ_N+\op(1).
\]
The smaller model minimizes the same quadratic subject to \(R\alpha=0\). The
ordinary equality-constrained quadratic-program identity gives
\[
  T_0-T_1
  =
  (R\hat\alpha_N)\T(RA^{-1}R\T)^{-1}(R\hat\alpha_N)+\op(1).
\]
Since
\[
  \hat\alpha_N\Rightarrow A^{-1}\Delta\T VZ,\qquad
  \Var(\hat\alpha)=A^{-1}BA^{-1},
\]
the covariance of \(R\hat\alpha\) is \(S\). The quadratic form
\((R\hat\alpha)\T C^{-1}(R\hat\alpha)\) therefore has the weighted chi-square
law in \eqref{eq:spectrum}, by simultaneous whitening of the pair \((C,S)\).
\end{proof}

\begin{corollary}[Normal-theory collapse]
\label{cor:nt}
If \(\Gamma=V^{-1}\), then \(B=A\), \(S=C\), and every nonzero eigenvalue in
\eqref{eq:spectrum} is one. Hence \(T_0-T_1\Rightarrow\chi^2_m\).
\end{corollary}

\begin{proof}
Substitution gives \(B=\Delta\T VV^{-1}V\Delta=A\), so
\[
  S=RA^{-1}AA^{-1}R\T=RA^{-1}R\T=C.
\]
\end{proof}

\section{FIML specialization}
\label{sec:fiml}

For raw incomplete continuous data, the saturated observed-data normal H1 model
has eta vector
\[
  \eta=(\mu,\vech\Sigma)
\]
per independent group or block. Let \(\hat\eta\) be the saturated EM/FIML
estimator. With observed-data normal log likelihood \(\ell_i(\eta)\), define
\[
  H_\eta
  =
  -\E\left\{\frac{\partial^2\ell_i(\eta_0)}
                  {\partial\eta\,\partial\eta\T}\right\},
  \qquad
  J_\eta
  =
  \E\left[
    \frac{\partial\ell_i(\eta_0)}{\partial\eta}
    \frac{\partial\ell_i(\eta_0)}{\partial\eta\T}
  \right].
\]
The first-stage sandwich covariance is
\[
  \Gamma_{\mathrm{mis}}=H_\eta^{-1}J_\eta H_\eta^{-1}.
\]
The saturated likelihood expansion around \(\hat\eta\) gives, for local
\(\eta=\eta_0+N^{-1/2}h\),
\[
  2\{\ell_N(\hat\eta)-\ell_N(\eta)\}
  =
  (Z_N-h)\T H_\eta (Z_N-h)+\op(1),
  \label{eq:fiml-local}
\]
so \eqref{eq:local-contrast} holds with \(V=H_\eta\) and
\(\Gamma=\Gamma_{\mathrm{mis}}\). If the larger SEM has local Jacobian
\[
  \Delta_1=\frac{\partial\eta_1(\theta)}{\partial\alpha\T}
\]
after equality-constraint tangent projection, and the smaller SEM is the local
restriction \(R\alpha=0\), Proposition~\ref{prop:restriction} gives
\[
  A_1=\Delta_1\T H_\eta\Delta_1,
  \qquad
  B_1=\Delta_1\T H_\eta\Gamma_{\mathrm{mis}}H_\eta\Delta_1,
\]
\[
  C=RA_1^{-1}R\T,
  \qquad
  S=RA_1^{-1}B_1A_1^{-1}R\T,
\]
and the Satorra--2000 mean scaling
\[
  \hat c=\frac{1}{m}\tr(C^{-1}S).
\]

This is exactly the route implemented by magmaan's FIML nested restriction map:
the code constructs \(V\) from saturated EM information, uses
\(\widehat\Gamma_{\mathrm{mis}}\) from saturated score meat, builds
\(\Delta_1\) in the same eta layout, obtains \(R\) from the H1-to-H0 local
restriction, and solves the generalized eigenproblem for the pair \((S,C)\).

\section{Why the estimated metric does not enter at this null}
\label{sec:estimated}

The FIML implementation uses finite-sample plug-in matrices, not fixed
population matrices. This does not change the first-order nested null law.
Let \(\hat V=V+\op(1)\), and let \(r_M(\hat\theta_M)\) be the fitted eta-space
residual for \(M\). If \(M_0\) is true in the saturated eta target, both
\[
  r_0(\hat\theta_0)=\Op(N^{-1/2}),\qquad
  r_1(\hat\theta_1)=\Op(N^{-1/2}).
\]
Therefore
\[
  N r_M(\hat\theta_M)\T(\hat V-V)r_M(\hat\theta_M)
  =
  \op(1).
\]
The same argument covers a smooth two-stage weight \(\hat W=W+\op(1)\) under
ML2S. Weight-estimation derivative terms are multiplied by the pseudo-true
residual. At the true nested null that residual is zero, so they vanish from
the first-order reference law.

The finite-sample scaling convention also cancels in the generalized
eigenproblem. magmaan stores the saturated information in summed
log-likelihood units, \(V_N\approx N V\), and the sandwich covariance in
\(1/N\) units, \(\Gamma_N\approx \Gamma/N\). Then
\[
  A_{1,N}=N A_1,\qquad B_{1,N}=N B_1,\qquad
  C_N=N^{-1}C,\qquad S_N=N^{-1}S,
\]
so the eigenvalues of \((S_N,C_N)\) equal those of \((S,C)\).

\section{ML2S specialization}
\label{sec:ml2s}

ML2S uses the same first-stage saturated EM estimator \(\hat\eta\), then fits
the structural model by a moment-quadratic criterion with a stage-two weight
\(W\). If \(W\) is treated as the local metric in
\eqref{eq:local-contrast}, Proposition~\ref{prop:restriction} applies with
\[
  A_1=\Delta_1\T W\Delta_1,
  \qquad
  B_1=\Delta_1\T W\Gamma_{\mathrm{mis}}W\Delta_1.
\]
For the normal-theory two-stage weight \(W=\Gamma_{\mathrm{NT}}^{-1}\) and
normal complete-data first-stage law, Corollary~\ref{cor:nt} again yields
unit eigenvalues. For empirical robust ML2S, \(\Gamma_{\mathrm{mis}}\) is the
saturated EM sandwich and the eigenvalues are the robust nested reference
spectrum.

When the stage-two weight is estimated from the same first-stage eta, its
derivative is still asymptotically irrelevant at the true nested null by the
argument in Section~\ref{sec:estimated}. It is not irrelevant when the nested
restriction is tested inside a misspecified larger model.

\section{Boundary of the derivation}
\label{sec:boundary}

The derivation above assumes the smaller model contains the first-stage target:
\[
  \eta_0\in M_0\subset M_1.
\]
This is the usual calibration null for a nested difference test and the regime
in which size is defined when the data-generating model satisfies the smaller
model. It is data-robust in the Satorra--Bentler/Yuan--Bentler sense: the
sandwich covariance may differ from the normal-theory covariance because the
observed data are nonnormal or incomplete, but the model restriction being
tested is still evaluated under a correctly specified smaller model.

A different null is possible: the larger and smaller profiled criteria have the
same pseudo-true optimum because the restriction is true inside a misspecified
larger model, but the common fitted residual \(r_*\) is nonzero. Then the first
derivative of the profiled difference still cancels, but the second derivative
is not the residual-projector covariance above. It is the profile Hessian
difference
\[
  Z\T(Q_0-Q_1)Z,
\]
where \(Q_M\) contains both model curvature terms and, when the metric is
estimated, derivatives of the metric map. This is the regime isolated in the
relative-fit and nested-profile note. The eta-space restriction map is recovered
from that profile-Hessian law only in special cases: fixed metric, zero
pseudo-true residual, or cancellations that make the profile Hessian equal the
ordinary Satorra projector.

\section{Implications for the lavaan parity gap}
\label{sec:parity}

The missing-data divergence from
\texttt{lavTestLRT(method = "satorra.2000")} is therefore not, by itself, an
oracle defect or a magmaan defect. It says that lavaan and magmaan are using
different missing-data conventions for the difference-test scaling quantity.
The eta-space route has a first-principles null derivation for saturated FIML
and ML2S under the true nested null. It is the robust method under study.

Two additional claims would require separate work.
\begin{enumerate}
\item Lavaan compatibility: derive and expose lavaan's missing-data finite-sample
  convention as a separate route, then gate that route against lavaan.
\item Misspecified nested robustness: implement or approximate the full
  profile-Hessian law for \(r_*\neq0\), including estimated-metric derivative
  terms, then calibrate it by simulation.
\end{enumerate}

\begin{thebibliography}{9}

\bibitem{satorra1989}
Satorra, A. (1989).
Alternative test criteria in covariance structure analysis: A unified approach.
\emph{Psychometrika}, 54(1), 131--151.
\href{https://doi.org/10.1007/BF02294453}{doi:10.1007/BF02294453}.

\bibitem{satorra2000}
Satorra, A. (2000).
Scaled and adjusted restricted tests in multi-sample analysis of moment
structures. In R. D. H. Heijmans, D. S. G. Pollock, \& A. Satorra (Eds.),
\emph{Innovations in Multivariate Statistical Analysis} (pp. 233--247).
Kluwer.
\href{https://doi.org/10.1007/978-1-4615-4603-0_17}{doi:10.1007/978-1-4615-4603-0\_17}.

\bibitem{satorrabentler2001}
Satorra, A., \& Bentler, P. M. (2001).
A scaled difference chi-square test statistic for moment structure analysis.
\emph{Psychometrika}, 66(4), 507--514.
\href{https://doi.org/10.1007/BF02296192}{doi:10.1007/BF02296192}.

\bibitem{savaleibentler2009}
Savalei, V., \& Bentler, P. M. (2009).
A two-stage approach to missing data: Theory and application to auxiliary
variables.
\emph{Structural Equation Modeling}, 16(3), 477--497.
\href{https://doi.org/10.1080/10705510903008238}{doi:10.1080/10705510903008238}.

\bibitem{white1982}
White, H. (1982).
Maximum likelihood estimation of misspecified models.
\emph{Econometrica}, 50(1), 1--25.
\href{https://doi.org/10.2307/1912526}{doi:10.2307/1912526}.

\bibitem{yuanbentler2000}
Yuan, K.-H., \& Bentler, P. M. (2000).
Three likelihood-based methods for mean and covariance structure analysis with
nonnormal missing data.
\emph{Sociological Methodology}, 30, 165--200.
\href{https://doi.org/10.1111/0081-1750.00078}{doi:10.1111/0081-1750.00078}.

\end{thebibliography}

\end{document}
